\section{Technische toelichting bij de toegankelijke hoofdtekst}\label{app:technical-nl}
De hoofdtekst gebruikt gewone taal en correcte aantallen. Hieronder staan de verplaatste technische termen, instellingen en exacte scores. De uitgewerkte vergelijking betreft RQ3; de definitieve methoden en resultaten voor RQ1/RQ2 zijn nog niet aangeleverd.

\subsection{Beeldselectie en validatie}
De bronverzameling bevat 78.000 A--Z-beelden. De historische featurecache accepteert 60.251 beelden met precies \mbox{\'e\'en} hand: 77,24\%. De classificatiescores zijn conditioneel op deze selectie. De meegeleverde Kaggle-testbeelden dupliceren trainingsbeelden en zijn uitgesloten. De historische detectiedrempels zijn niet bewaard; overeenstemming met nieuw gedetecteerde beelden mag niet worden verondersteld.

Er zijn vijf vaste validatiefolds over zes bestandsnummerblokken, geen geverifieerde personen. Per fold wordt op vier delen getraind en op het overgebleven deel voorspeld. Elk geaccepteerd beeld krijgt \mbox{\'e\'en} out-of-fold voorspelling per variant. Afgeleide beelden blijven bij hun bron. Imputatie en schaling van toegevoegde hoeken worden uitsluitend op het trainingsdeel geleerd. Bij een leercurve moeten deze bewerkingen opnieuw binnen iedere trainingssubset worden geleerd.

\subsection{Representaties en leermethoden}
De Shin-representatie heeft 820 waarden: 190 genormaliseerde afstanden en 630 richtingswaarden. Het volledige toevoegingspakket bevat 70 categorische inputs, 26 overeenkomstscores met het werkboek en 15 gewrichtsmetingen. De overeenkomstscores worden voor alle kandidaatletters berekend, zonder het ware label als invoer. De 15 metingen bestaan uit 14 interne gewrichtshoeken en een getekende duimafwijking. Het zijn handpuntmetingen, geen klinisch gevalideerde gewrichtshoeken.

De wereldvariant gebruikt MediaPipes geschatte 3D-wereldpunten. De controle gebruikt exact dezelfde genormaliseerde beeldpunten als Shin, inclusief geschatte diepte. Die controle voegt 15 waarden toe (835 inputs), of vervangt de richtingen door hoeken (205 inputs). De laatstgenoemde vervanging is een afzonderlijke proef. Zij haalt 94,511\% nauwkeurigheid met SVM; dit is geen resultaat van het toevoegen van hoeken.

Het vaste SVM-profiel gebruikt een RBF-kernel, $C=10$, \texttt{gamma=scale} en gebalanceerde klassen. LightGBM gebruikt 100 bomen, 31 bladeren, leersnelheid 0,1 en minimaal 20 voorbeelden per blad. Hoeken worden met de trainingsmediaan aangevuld en gestandaardiseerd. De resultaten zeggen niet welk algoritme na een even uitgebreide optimalisatie het beste zou zijn. Verdere formules, controles en alle letters staan in bijlagen A--K.

\subsection{Scores en exacte waarden}
Nauwkeurigheid is het aandeel correcte voorspellingen. Voor iedere letter worden true positives (TP), false positives (FP) en false negatives (FN) geteld:
\[
P=\frac{TP}{TP+FP},\qquad
R=\frac{TP}{TP+FN},\qquad
F1=\frac{2PR}{P+R}.
\]
Hier is $P$ precisie en $R$ recall. Niet-gedefinieerde scores worden nul. Macro-precisie en macro-F1 zijn de ongewogen gemiddelden over A--Z. Gepoolde scores worden uit alle out-of-fold voorspellingen samen berekend, niet als ongewogen gemiddelde van de vijf folds. Bij de spiegelstudie worden niet-geaccepteerde detecties bovendien als gemiste pogingen behouden; die populatie is anders dan de 60.251 geaccepteerde beelden.

\begin{table}[h]\centering\small
\caption{Exact dezelfde SVM-resultaten als de aantallentabel in de hoofdtekst, nu met percentages en macro-F1.}
\begin{tabular}{lrr}\toprule
Variant & Nauwkeurigheid (\%) & Macro-F1\\\midrule
\tableinput compact_results_nl.tex
\bottomrule\end{tabular}
\end{table}

\clearpage
\subsection{Aanvullende uitkomsten en nog open bijdragen}
De wereldvariant herstelt 143 fouten en verliest 76 goede voorspellingen (netto +67). Dezelfde-co\"ordinatenvariant herstelt 89 en verliest 85 (netto +4). Bij LightGBM verandert het volledige pakket van 97,097\% naar 97,164\% nauwkeurigheid. Deze verschillen zijn verkennende uitkomsten op hergebruikte folds, geen bevestigende significantietoets.

Het eerdere RQ1-concept vergelijkt een tweede detectiepoging met $3\times3$-dilatatie, erosie en een Canny-overlay. Dilatatie neemt lokaal een maximale pixelwaarde; dit is geen verwijdering. Definitieve instellingen, herstelaantallen en visuele kwaliteitscontrole blijven TODO. Herstelpercentages gebruiken de oorspronkelijke gemiste beelden als noemer en moeten unieke beelden onderscheiden van overlappende filterresultaten.

Voor RQ2 moeten aantallen per klasse, geneste trainingssubsets, vaste toets-ID's, seeds en herhalingen worden vastgelegd. Het aantal unieke bronbeelden is de variabele, niet het aantal augmentatiekopie\"en. Detector, representatie, classifier en toetsverzameling blijven gelijk; variatie tussen herhalingen moet worden gerapporteerd. Deze opzet is een voorstel, geen claim over reeds uitgevoerde proeven van de medestudent.

De spiegelstudie (bijlage L) verandert SVM-macro-F1 op de vaste halfgespiegelde presentatie van 0,53849 naar 0,84560, met detectiefalen meegeteld. Op originele beelden gaan netto 26 goede voorspellingen verloren. De losse G/H-proef (bijlage M) gebruikt alleen $180^\circ$ camera-$y$-rotatie van trainingsrichtingen, $(d_x,d_y,d_z)\mapsto(-d_x,d_y,-d_z)$, met behoud van brongewicht en oorspronkelijke preprocessing. Er is geen spiegeltraining, camerakalibratie of routering. Macro-F1 verandert van 0,97524391 naar 0,97523305: vier hersteld en vijf verloren. Een geometrische draaiing genereert geen werkelijk nieuw RGB-aanzicht of bijbehorende detectorfouten.
